% Source: https://artofproblemsolving.com/community/c3h1256108_solving_similar_triangles_relating_segments
% Stable identity: aops-topic-1256108-q01
Triangles $CBD$ and $CAB$ share the angle at $C$, and $\angle DBC=\angle BAC$, so they are similar (in that order).

\begin{asy}
usepackage("amsmath");
size(4.5cm);
defaultpen(fontsize(10pt));
pair A=(0,0),B=(4,0); pair D=4*dir(40);
pair Cc=extension(A,A+dir(40),B,B+dir(70));
draw(A--B--Cc--cycle,linewidth(1)); draw(B--D,blue); dot(D);
label("$A$",A,SW);label("$B$",B,SE);label("$C$",Cc,N);label("$D$",D,NW);
label("\mbox{}",(min(currentpicture,user=true)+max(currentpicture,user=true))/2,black); // keeps the text after the figure black when LaTeX typesets the labels inline
\end{asy}

Hence
\[\dfrac{CB}{CA}=\dfrac{CD}{CB}=\dfrac{BD}{AB}.\]
Write $a=AD=AB$ and $d=DC$, so $CA=a+d$. Then $CB^2=CD\cdot CA=d(a+d)$, so $CB=\sqrt{d(a+d)}$, and $BD=\dfrac{AB\cdot CB}{CA}=\dfrac{a\,CB}{a+d}$. Therefore
\[BC-BD=CB\left(1-\dfrac{a}{a+d}\right)=\dfrac{d\sqrt{d(a+d)}}{a+d}=\boxed{DC\sqrt{\dfrac{DC}{AD+DC}}}.\]
